\section{Formalization plan and reproducibility}
\label{sec:validation}

\subsection{Lean dependency plan}

The new algebra can be formalized independently of the full theta-lift
construction. A natural sequence is:
\begin{enumerate}
\item For representations spanning the full matrix algebra, prove
absolute irreducibility of the residue representations and the
one-dimensional bound for the residue intertwiner space.
\item Identify the uniformizer torsion of the finite-level Hom module
with that residue Hom. Use the structure theorem for finite-length
modules over a DVR to deduce cyclicity.
\item Combine cyclicity with
\texttt{smul\_apply\_eq\_zero\_of\_mem\_conjDualHom}.
This already proves the sharp scalar estimate.
\item Formalize trace rigidity over \(O/\pi^n\) using the perfect trace
pairing, matrix units, and the innerness of matrix-algebra
automorphisms over a local ring.
\item Deduce exact trace depth, finite generator tests, and the
conjugate-dual identification. Then formalize the Ramanujan example
relative to the classical representation and image theorems.
\end{enumerate}
This is a dependency plan, not a claim that these Lean files have been
compiled. No new \texttt{sorry}-containing theorem is packaged as a
completed formal proof.

If \(e=v_\pi(\ell)\), a raw length \(a\) is \(a/e\) under the source's
normalized length. The main scalar corollary must preserve that
normalization. In the worked \(\Z_\ell\) example \(e=1\).

\subsection{What the exact checkers establish}

\begin{center}
\small
\begin{tabular}{p{0.33\linewidth}p{0.59\linewidth}}
\toprule
Checker & Verified content\\
\midrule
\path{verify_trace_depth.py} &
Unimodular matrix and trace-Gram bases; Smith factors;
exhaustive small quotient kernels; dyadic transpose scalars;
the two-prime Ramanujan certificate and finite cyclotomic-unit tests\\[3pt]
\path{verify_transfer_extensions.py} &
Finite trace-test identities; exact nonmaximal image orders and
extra torsion; the common denominator ideal example\\[3pt]
\path{euler_factor_check.py} &
Rank-two identities on small exact matrices and the
rank-three counterexample\\[3pt]
\path{verify_density_gap_independent.py} &
All subsets through 12, the exact prefix loss \(1/156\),
and a rational lower certificate exceeding \(0.000404\)\\[3pt]
\path{upstream_audit.py} &
Pinned source identity, static token scan, source-file counts,
and hashes of inspected files\\
\bottomrule
\end{tabular}
\end{center}

The first two checkers use SymPy for integer Smith normal forms.
The other three use the Python standard library.
The Ramanujan unit checks run over all units in the specified finite
quotients; the theorem for every Tate weight and every good prime is
proved by the group-theoretic and valuation argument in the report.
Likewise, no finite kernel enumeration proves a general representation
theorem or certifies a Galois realization.

Two larger inherited density checkers were inspected and replayed.
Their original algorithms, results, and exact provenance are included
under \path{evidence/prior_density/}. The independent checker in this
package is a separate implementation of the simpler common gap.

\subsection{Package contents and rebuilding}

The archive contains this report's complete TeX source and compiled PDF,
the five checkers, saved JSON certificates, the 46-row prior-report
ledger, selected upstream source files with license and attribution,
the mathematical review notes, and a formalization plan.
\path{README.txt} gives exact commands.
\path{MANIFEST.sha256} records the packaged file hashes.
The selected source files are evidence snapshots, not a standalone
copy of all upstream transitive dependencies.

The PDF was compiled from the supplied TeX and visually inspected after
rendering. The checks described here are paper-proof review, exact
computation, source inspection, and document validation. They are
distinct from the independent Lean kernel replays cited in
Section~\ref{sec:scope}.
